\documentclass[a4paper,10pt]{article}
\usepackage[UTF8]{ctex}
\usepackage[margin=1.2cm]{geometry}
\usepackage{titlesec}
\usepackage{enumitem}
\usepackage{xcolor}
\usepackage{hyperref}

\definecolor{sectioncolor}{RGB}{23, 63, 122}
\titleformat{\section}{\large\bfseries\color{sectioncolor}}{}{0em}{}[\titlerule]
\titlespacing{\section}{0pt}{8pt}{5pt}
\setlist[itemize]{leftmargin=*,itemsep=1pt,parsep=1pt,topsep=2pt}
\hypersetup{colorlinks=true,urlcolor=sectioncolor}
\pagestyle{empty}

\begin{document}

\begin{center}
    {\huge\bfseries 李 \ 上 \ 一} \\
    \vspace{6pt}
    山东大学 · 计算机科学与技术 \quad | \quad 政治面貌：共青团员 \\
    邮箱：\href{mailto:202300130120@mail.sdu.edu.cn}{202300130120@mail.sdu.edu.cn}
    \quad | \quad
    GitHub：\href{https://github.com/SDLSY}{github.com/SDLSY}
\end{center}

\section{教育背景}
\textbf{山东大学 \ (Shandong University)} \hfill \textbf{2023.09 -- 2027.07} \\
\textbf{专业：}计算机科学与技术 \hfill
\textbf{GPA：}91.28 / 100 \hfill
\textbf{专业成绩排名：}12 / 228 \\
\textbf{研究兴趣：}可穿戴计算、移动感知、健康计算、边缘智能、端侧 AI \\
\textbf{主修课程：}计算机网络(98)、离散数学(97)、高等数学 I/II(97/94)、数据库系统(95)、线性代数(94)、计算机组成原理(93)、概率论(92)、电路(100) \\
\textbf{英语水平：}CET-4 551 \quad | \quad CET-6 489 \quad | \quad 英语演讲与辩论 97

\section{科研与项目经历}
\textbf{长庚环 VesperO：智能戒指端云协同健康辅助系统} \hfill \textbf{2024.11 -- 至今} \\
\textit{全国大学生软件创新大赛全国三等奖；团队项目，除视频制作外，核心开发、调试集成与材料整理主要由本人完成。}
\begin{itemize}
    \item 面向睡眠恢复与居家健康管理场景，构建智能戒指、Android App、云端服务、AI 问诊、报告理解与后台管理组成的端云协同系统。
    \item 负责 Android 多模块应用与核心链路开发，完成 BLE 设备接入、实时数据展示、本地数据管理、云端接口对接和 AI 辅助功能调试。
    \item 接入真实智能戒指硬件，结合真实采集与演示数据，覆盖 PPG、心率、血氧、体温、三轴加速度和三轴陀螺仪等多模态数据。
    \item 搭建本地异常检测、PPG/HRV/体温分析、睡眠分期服务和健康建议生成流程，形成从数据采集到反馈展示的业务闭环。
\end{itemize}

\textbf{基于戒指式可穿戴设备的盲文触摸信号识别} \hfill \textbf{2025.10 -- 至今} \\
\textit{在老师指导下开展；面向智能戒指 IMU 信号的盲文触摸动作识别与端侧实时识别探索。}
\begin{itemize}
    \item 使用 Hi90B 智能指环采集 6 轴 IMU 数据，构建 26 个英文盲文字母触摸动作的单被试初步实验数据集，每类不少于 100 组。
    \item 搭建 Android 采集、CSV 清洗、信号分段、特征提取和分类评估流水线，比较手工时序特征、MiniROCKET、MultiRocketHydra 与小波去噪等方法。
    \item 在严格时序划分下将识别准确率由约 81\% 提升至 87.2\%，并分析单会话采集带来的 session bias 与跨用户泛化限制。
\end{itemize}

\textbf{可视计算暑期课堂：基于目标跟踪的增强现实 (AR) 系统} \hfill \textbf{2024.07 -- 2024.08}
\begin{itemize}
    \item 使用 ORB + RANSAC 实现特征点稳健匹配，完成相机内参标定、3D 空间坐标变换和移动端实时渲染，达到 30FPS 虚实融合效果。
\end{itemize}

\textbf{OceanBase / MiniOB 数据库大赛性能调优} \hfill \textbf{2024.09 -- 2025.01}
\begin{itemize}
    \item 参与 MiniOB 存储与查询优化，接触 B+ 树索引、LRU-K 缓冲池和并发控制机制，获 OceanBase 大赛性能调优优胜。
\end{itemize}

\section{竞赛与荣誉}
\begin{itemize}
    \item \textbf{学科竞赛：}全国大学生软件创新大赛全国三等奖、华东赛区二等奖；蓝桥杯 C++ 省一等奖；全国大学生数学建模竞赛省级二等奖；OceanBase 大赛性能调优优胜。
    \item \textbf{奖学金与荣誉：}一等学业奖学金(2024)、二等学业奖学金(2025)、校级三好学生(连续两年)、山东省优秀社会实践奖。
\end{itemize}

\section{专业技能}
\begin{itemize}
    \item \textbf{编程与开发：}熟悉 C/C++、Python、Java/Kotlin Android 开发，具备 Linux、Git、移动端工程和端云协同开发经验。
    \item \textbf{智能感知与信号处理：}熟悉 BLE 通信、传感器数据采集、时序信号预处理、小波变换、卡尔曼滤波和基础机器学习实验流程。
    \item \textbf{系统与工程：}了解数据库索引、缓存管理和并发控制，具备 Room/SQLite、Supabase、Next.js 云端服务和 ONNX 推理服务集成经验。
\end{itemize}

\end{document}
